% Pista ana-25s-q1 · Anàlisi
% Procedència: PAU Catalunya 2025 · Convocatòria de setembre · Sèrie 3 · Exercici 1
\documentclass[11pt,a4paper]{article}
\usepackage{pista}
\pistainfo{Catalunya 2025}{Convocatòria de setembre}{Sèrie 3}

\begin{document}

\section*{\textcolor{blauPista}{Exercici 1}}

\noindent Una família vol comprar un terreny per a fer-s'hi una casa envoltada de penya-segats amb vistes al mar. En aquella zona de la costa, els penya-segats segueixen les rectes $y = 0$ i $y = 3x$. A més, la família vol que el terreny sigui triangular i que el tercer costat del triangle passi pel punt $P = (1, 1)$, tal com es pot veure en la figura.

\subsection*{\textit{a)} \normalfont Plantegeu l'equació de la recta $r$ que defineix el tercer costat del triangle en funció del seu pendent $m$, i comproveu que l'àrea del terreny ve donada per
\[
A(m) = \dfrac{3}{2} \cdot \dfrac{m^{2} - 2m + 1}{m^{2} - 3m}.
\]
\hfill\textnormal{[1~punt]}}

\pista{L'equació de la recta $r$ et ve donada pel punt $P$ i el pendent $m$ (forma punt-pendent). Un cop tinguis $r$, troba els dos punts de tall d'aquesta recta amb les rectes dels penya-segats, $y = 0$ i $y = 3x$. Amb això, obtindràs dos vèrtexs del triangle, i el tercer és el $(0,0)$. Finalment, aplica la fórmula de l'àrea d'un triangle a partir dels vèrtexs, i simplifica el resultat.}

\subsection*{\textit{b)} \normalfont Calculeu el valor de $m$ que fa que l'àrea d'aquest terreny (i, per tant, el seu preu) sigui mínima. Quin és el valor d'aquesta àrea? \hfill\textnormal{[1,5~punts]}}

\pista{És un problema d'optimització: deriva $A(m)$ i resol $A'(m) = 0$. Obtindràs un o més valors candidats: descarta els que no tinguin sentit geomètric (per exemple, els que facin degenerar el triangle o els que corresponguin a una recta paral·lela a un dels penya-segats). Comprova que el valor triat és efectivament un mínim i substitueix-lo a $A(m)$. D'aquesta manera, trobaràs l'àrea mínima.}

\end{document}
